\documentclass[a4paper,11pt]{article}

\usepackage{fontspec}
\usepackage{newunicodechar}
\newunicodechar{π}{\ensuremath{\pi}}
\usepackage[a4paper,margin=1in]{geometry}
\usepackage{hyperref}
\hypersetup{colorlinks=true,linkcolor=blue,citecolor=blue,urlcolor=blue}
\usepackage{enumitem}
\usepackage{booktabs}
\usepackage{tabularx}
\usepackage{array}
\usepackage{float}
\newcolumntype{Y}{>{\raggedright\arraybackslash}X}
\usepackage[numbers]{natbib}
\usepackage{url}
\emergencystretch=2em
\sloppy

\title{Evaluating Retrieval-Augmented Generation Systems}
\author{}
\date{\today}

\setlist[itemize]{noitemsep,topsep=0.25em,leftmargin=*}
\setcounter{tocdepth}{2}

\begin{document}
\maketitle

\begin{abstract}


Recent research on the evaluation of retrieval-augmented generation systems now spans multiple methodological families, task settings, and evaluation regimes, yet those threads are still often compared under mismatched protocols \citep{Sun2026Robust,Kim2024Learning,Salemi2024Evaluating,Taiwo2026Evaluating,Kovalskii2026Domain,Tian2026Predicting}. Across work on retrieval and generation failure decomposition, measurement validity and evaluator design, and robustness, transfer, and deployment, the same question keeps returning: which reported gains survive changes in setup, data, and evaluation practice \citep{Caruzzo2026Deceptive,Ding2026Attribution,Murugaraj2025Ragvue,Nematov2025Source,SaadFalcon2023Ares,Upadhyay2026Overview}. Drawing on a candidate pool of 320 records and a 48-paper core set, the paper tracks how current approaches trade off effectiveness, generality, reproducibility, and deployment realism across representative settings \citep{Yao2026Authoritybench,An2026Fresco,Li2026Mrag,Guder2026Rage,Brehme2025Retrieval,Packowski2024Optimizing}. A central finding is that many reported gains depend as much on evaluation design and experimental assumptions as on the nominal methodology itself \citep{Muthyam2026Systems,Feng2026Rethinking,Su2025Dynamic,Bang2026Feedback,Cahoon2025Optimizing,Javadi2025Contradictions}. The conclusion identifies where current evidence is strongest, where protocol mismatch still obscures comparison, and which open questions should structure future work \citep{Wampler2025Engineering,Zhou2024Trustworthiness}.
\end{abstract}

\section{Introduction}


Recent work on the evaluation of retrieval-augmented generation systems asks systems to operate across heterogeneous tasks, operating contexts, and deployment conditions that do not line up as neatly as the shared labels imply. Closely related results frequently shift the task definition, information access, or deployment constraint that gives a reported gain its meaning. Any synthesis that ignores those differences risks turning protocol mismatch into a false methodological ranking \citep{Sun2026Robust,Kim2024Learning,Salemi2024Evaluating,Taiwo2026Evaluating}.

Across that literature, the same comparison problem keeps returning in different guises: retrieval diagnostics, protocol design, robustness under shift. They are more stable than method names alone because they keep the comparison attached to the mechanism, protocol, and limitation evidence that travel with each claim. Keeping those questions in one frame lets retrieval diagnostics, protocol design, robustness under shift remain interpretable within a single argument instead of scattering into disconnected local summaries \citep{Kovalskii2026Domain,Tian2026Predicting,Caruzzo2026Deceptive,Ding2026Attribution}.

The survey draws on a candidate pool of 320 records and narrows that pool to a 48-paper core set under abstract-level evidence. The synthesis uses recent years as its retrieval window, so later comparisons inherit a visible scope boundary instead of an implied one. That boundary matters because later claims should be read against the evidence window that produced them, not against an implied exhaustive history \citep{Murugaraj2025Ragvue,Nematov2025Source,SaadFalcon2023Ares,Upadhyay2026Overview}. Within that frame, stronger claims are reserved for papers that make task, metric, and constraint choices explicit. Where protocols diverge too sharply, the paper uses the citations to delimit the comparison rather than to manufacture a false consensus. The practical consequence is a stricter distinction between stable evidence and suggestive but not yet decisive comparisons \citep{Yao2026Authoritybench,An2026Fresco,Li2026Mrag,Guder2026Rage}.
\section{Related Work}


The field around the evaluation of retrieval-augmented generation systems already includes several useful reviews that clarify local vocabularies, benchmark choices, and subarea boundaries. What is still unstable is less the existence of prior overviews than the fact that many of them combine orientation, benchmarking, and application framing in ways that make cross-paper comparison drift. The key question is therefore not whether prior surveys exist, but which comparisons they actually stabilize once their assumptions are unpacked \citep{Sun2026Robust,Kim2024Learning,Salemi2024Evaluating,Taiwo2026Evaluating}.

One useful strand of prior work around retrieval and generation failure decomposition earns its value by stabilizing what counts as a like-for-like comparison inside that slice. Those reviews sharpen local baselines and task definitions, but they rarely keep retrieval diagnostics aligned once the evidence starts travelling across adjacent chapters. Their value is therefore diagnostic: they show where the local vocabulary is stable and where the synthesis still has to rebind the evidence itself \citep{Kovalskii2026Domain,Tian2026Predicting,Caruzzo2026Deceptive,Ding2026Attribution}.

Survey work on measurement validity and evaluator design matters less as a ready-made outline than as evidence about which assumptions that slice already treats as fixed. That literature is useful for orientation, yet it usually inherits the assumptions of its local slice rather than re-reading protocol design against neighboring evidence. That makes them essential orientation, but not a substitute for the cross-chapter comparison basis the body still has to build \citep{Murugaraj2025Ragvue,Nematov2025Source,SaadFalcon2023Ares,Upadhyay2026Overview}. Benchmark- and evaluation-centered papers contribute a different kind of related work. They expose metric choice, task coverage, and reproducibility as substantive methodological decisions, which is often exactly where apparent disagreements between methods become interpretable. In practice, those papers are useful less as another category label and more as calibration for what a fair comparison can mean \citep{Yao2026Authoritybench,An2026Fresco,Li2026Mrag,Guder2026Rage}.

The unresolved problem is not naming one more family of systems. Existing reviews often settle one of those dimensions well while letting the others recede into background assumptions. Without that joint frame, the literature drifts back toward inventories even when the cited evidence is rich enough to support stronger synthesis \citep{Brehme2025Retrieval,Packowski2024Optimizing,Muthyam2026Systems,Feng2026Rethinking}. The value of prior surveys here is less as inherited structure than as evidence about what the field already treats as settled or unstable. Related work stays useful here only when it remains a positioning layer instead of turning into a surrogate table of contents. The benefit is a related-work section that positions the reader before the comparisons begin instead of pre-writing the argument for them \citep{Su2025Dynamic,Bang2026Feedback,Cahoon2025Optimizing,Javadi2025Contradictions}.
\section{Evaluation Targets \& Failure Decomposition}


The core issue in Evaluation Targets \& Failure Decomposition is not breadth alone, but how retrieval diagnostics and faithfulness and attribution reshape what counts as a fair comparison \citep{Sun2026Robust,Kim2024Learning,Salemi2024Evaluating,Taiwo2026Evaluating}. Retrieval quality and evidence coverage and Generation faithfulness and source attribution each expose that pressure from a different angle, showing why retrieval diagnostics and faithfulness and attribution cannot be treated as interchangeable background choices \citep{Kovalskii2026Domain,Tian2026Predicting,Caruzzo2026Deceptive,Ding2026Attribution}.
\subsection{Retrieval quality and evidence coverage}


The literature on retrieval quality and evidence coverage becomes hard to compare when papers keep the same label but change the conditions that give their gains meaning. Retrieval evaluation should separate rank quality, evidence coverage, and downstream answer utility because each answers a different question \citep{Salemi2024Evaluating,Taiwo2026Evaluating,Sun2026Robust,Kim2024Learning}. This is the central pressure point: whether quality is measured as rank relevance, evidence coverage, or downstream answer utility, because each supports a different interpretation \citep{Salemi2024Evaluating,Taiwo2026Evaluating,Sun2026Robust,Kim2024Learning}. Evaluation of the retrieval model's performance based on query-document relevance labels shows a small correlation with the RAG system's downstream performance \citep{Salemi2024Evaluating}.

Papers grouped under retrieval diagnostic studies sharpen the contrast in retrieval quality and evidence coverage because changing those assumptions changes what the method is actually built to do. The findings show that structure-aware chunking yields higher overall retrieval effectiveness, particularly in top-K metrics, and incurs significantly lower computational costs than semantic or baseline strategies \citep{Taiwo2026Evaluating}. In contrast, a second strand, grouped here as retrieval utility and coverage optimization, shifts which assumptions the comparison has to keep fixed. Trade-off evidence from benchmark tasks does not travel cleanly across retrieval quality and evidence coverage without matching budget and cost model conditions \citep{Sun2026Robust}. The cited trade-off remains local to benchmark tasks, with its budget and cost model conditions shaping the comparison \citep{Kim2024Learning}.

Across this subsection, the meaningful split between retrieval diagnostic studies and retrieval utility and coverage optimization is which assumptions each route preserves once the setting starts to move \citep{Salemi2024Evaluating,Taiwo2026Evaluating,Sun2026Robust,Kim2024Learning}.

The hardest part of reading retrieval quality and evidence coverage is deciding which results survive a like-for-like comparison. Existing multi-path reasoning methods improve robustness by sampling multiple candidate answers and applying voting- or confidence-based selection, but they still face two limitations: diversity is often injected through uncontrollable decoding randomness, and answer evaluation is usually confined to a single query-induced evidence view \citep{Sun2026Robust}. A key challenge for improving RAG is to predict both the utility of retrieved documents -- quantified as the performance gain from using context over generation without context -- and the quality of the final answers in terms of correctness and relevance \citep{Tian2026Predicting}. Cross-paper judgment remains fragile until work on retrieval diagnostic studies and retrieval utility and coverage optimization exposes task, metric, and deployment constraints in comparable form \citep{Taiwo2026Evaluating,Salemi2024Evaluating}.
\subsection{Generation faithfulness and source attribution}


The central decision in generation faithfulness and source attribution is not which route sounds stronger in the abstract, but which assumptions remain stable once the setting changes. Generation evaluation should distinguish claim support, source-attribution correctness, citation quality, and answer utility rather than treating faithfulness as one score \citep{Ding2026Attribution,Nematov2025Source,SaadFalcon2023Ares,Murugaraj2025Ragvue}. The evidence is hardest to interpret here: claim-level support, attribution correctness, and answer utility can point in different directions, so evaluator choice remains part of the conclusion \citep{Ding2026Attribution,Nematov2025Source,SaadFalcon2023Ares,Murugaraj2025Ragvue}. In the construct with the most multi-dataset human-labeled coverage -- generated-answer attribution (AttributionBench's four source datasets, n = 1,610, with independent HAGRID, n = 2,150) -- none does: the per-dataset metric rankings invert (Kendall tau = -0.64, p = 0.031 on AttributedQA vs. LFQA), and an off-the-shelf NLI scorer that is best on short-claim AttributedQA (AUROC 0.90) collapses to AUROC 0.53 (chance) on long-form LFQA, where BERTScore wins (0.91); the flip is not a length or truncation artifact \citep{Ding2026Attribution}.

A recurring strand in generation faithfulness and source attribution keeps the comparison anchored in attribution and grounding diagnostics, where those assumptions directly shape what the system is optimizing. The primary obstacle is the substantial computational cost, where each utility function evaluation involves an expensive LLM call, resulting in direct monetary and time expenses \citep{Nematov2025Source}. Deceptive grounding (DG) describes a failure invisible to faithfulness, hallucination, and citation checks because every claim is sourced from a real document, about the wrong entity \citep{Caruzzo2026Deceptive}.

Evidence for attribution and grounding diagnostics only travels beyond one paper when those assumptions stay visible in the reporting contract. Attribution methods, such as Shapley values, are widely used to explain the importance of features or training data in traditional machine learning, their application to Large Language Models (LLMs), particularly within Retrieval-Augmented Generation (RAG) systems, is nascent and challenging \citep{Nematov2025Source}. One local evidential boundary still matters here because this instability has a concrete decision cost: a naive "best-on-average" rule for choosing an evaluator fails leave-one-dataset-out (mean held-out regret 0.172 AUROC, worse than fixing one scorer), so metric choice must be validated on the target dataset rather than learned from others \citep{Ding2026Attribution}.

In contrast, seen through automated faithfulness evaluators, the disagreement moves away from attribution and grounding diagnostics and toward a different set of assumptions. ARES is an Automated RAG Evaluation System for evaluating RAG systems along the dimensions of context relevance, answer faithfulness, and answer relevance \citep{SaadFalcon2023Ares}. Evaluating Retrieval-Augmented Generation (RAG) systems remains a challenging task: existing metrics often collapse heterogeneous behaviors into single scores and provide little insight into whether errors arise from retrieval, reasoning, or grounding \citep{Murugaraj2025Ragvue}.

Across generation faithfulness and source attribution, attribution and grounding diagnostics and automated faithfulness evaluators diverge less on headline performance than on the assumptions that make their results travel \citep{Ding2026Attribution,Nematov2025Source,SaadFalcon2023Ares,Murugaraj2025Ragvue}. The hardest part of reading generation faithfulness and source attribution is deciding which results survive a like-for-like comparison. A prompt-based LLM judge avoids the chance-level collapses the automatic scorers suffer (no LFQA collapse) but is not uniformly best, \textasciitilde{}100x costlier, and non-deterministic -- relocating, not removing, the validation burden \citep{Ding2026Attribution}. What is still missing is reporting that keeps task, metric, and deployment details visible when attribution and grounding diagnostics and automated faithfulness evaluators are compared \citep{Caruzzo2026Deceptive,Ding2026Attribution}.
\section{Evaluation Protocols \& Measurement Validity}


In Evaluation Protocols \& Measurement Validity, the central question is how protocol design and evaluator choice change the meaning of reported progress \citep{Yao2026Authoritybench,An2026Fresco,Li2026Mrag,Upadhyay2026Overview}. Dataset, task, and metric design and Automated evaluators and human validation each expose that pressure from a different angle, showing why protocol design and evaluator choice cannot be treated as interchangeable background choices \citep{Guder2026Rage,SaadFalcon2023Ares,An2026Fresco,Brehme2025Retrieval}.
\subsection{Dataset, task, and metric design}


The comparison in dataset, task, and metric design only holds when those assumptions remain visible instead of being absorbed into a generic benchmark story. Dataset and metric design determine which RAG failures are observable, so task coverage, reference assumptions, and aggregation rules must be reported together \citep{Guder2026Rage,SaadFalcon2023Ares,Yao2026Authoritybench,Li2026Mrag}. The comparison becomes unstable at this point: protocol breadth and construct validity do not automatically align: broader task coverage can make metric interpretation and reference assumptions less consistent \citep{Guder2026Rage,SaadFalcon2023Ares,Yao2026Authoritybench,Li2026Mrag}. Trade-off evidence from benchmark tasks does not travel cleanly across dataset, task, and metric design without matching budget and cost model conditions \citep{Upadhyay2026Overview}.

Papers grouped under evaluation frameworks and shared tracks sharpen the contrast in dataset, task, and metric design because changing those assumptions changes what the method is actually built to do. Across eight different knowledge-intensive tasks in KILT, SuperGLUE, and AIS, ARES accurately evaluates RAG systems while using only a few hundred human annotations during evaluation \citep{SaadFalcon2023Ares}. The evaluated approach leverages core techniques in LLM applications, including document chunking, vector databases, embedding models, and retrievers, to evaluate trade-offs among accuracy, efficiency, and scalability \citep{Guder2026Rage}.

In contrast, seen through domain and stress-test benchmarks, the disagreement moves away from evaluation frameworks and shared tracks and toward a different set of assumptions. Retrieval-Augmented Generation (RAG) has been swiftly adopted in scientific and clinical QA systems, a comprehensive evaluation benchmark in the medical domain is lacking \citep{Li2026Mrag}. Existing benchmarks predominantly evaluate re-rankers in static settings and do not adequately assess performance under evolving information -- a critical gap, as real-world systems often must choose among temporally different pieces of evidence \citep{An2026Fresco}.

Across dataset, task, and metric design, evaluation frameworks and shared tracks and domain and stress-test benchmarks diverge less on headline performance than on the assumptions that make their results travel \citep{Guder2026Rage,SaadFalcon2023Ares,Yao2026Authoritybench,Li2026Mrag}. The hardest part of reading dataset, task, and metric design is deciding which results survive a like-for-like comparison. A consistent failure mode across existing re-rankers: a strong bias toward older, semantically rich documents, even when they are factually obsolete \citep{An2026Fresco}. A sharper answer depends on knowing when gains survive changes in task, operating environment, and budget rather than only one benchmark \citep{Guder2026Rage,SaadFalcon2023Ares}.
\subsection{Automated evaluators and human validation}


The useful lens in automated evaluators and human validation is not a single architecture name, but the coupled choices that give the comparison its meaning. Automated evaluators are useful only to the extent that their scores remain calibrated against human judgments and the constructs they are intended to measure \citep{SaadFalcon2023Ares,Murugaraj2025Ragvue,Brehme2025Retrieval,Caruzzo2026Deceptive}. The evidence is hardest to interpret here: automated judging improves scale and repeatability, while human validation remains the reference point for unfamiliar failures and construct validity \citep{SaadFalcon2023Ares,Murugaraj2025Ragvue,Brehme2025Retrieval,Caruzzo2026Deceptive}. Common RAG benchmark evaluation techniques have not been useful for evaluating responses to novel user questions, so the evidence indicates a flexible, "human in the lead" approach is required \citep{Packowski2024Optimizing}.

Papers grouped under automated and model-based evaluators sharpen the contrast in automated evaluators and human validation because changing those assumptions changes what the method is actually built to do. This instability has a concrete decision cost: a naive "best-on-average" rule for choosing an evaluator fails leave-one-dataset-out (mean held-out regret 0.172 AUROC, worse than fixing one scorer), so metric choice must be validated on the target dataset rather than learned from others \citep{Ding2026Attribution}.

The evidence for automated and model-based evaluators in automated evaluators and human validation depends less on the headline number than on whether those assumptions remain comparable across papers. A prompt-based LLM judge avoids the chance-level collapses the automatic scorers suffer (no LFQA collapse) but is not uniformly best, \textasciitilde{}100x costlier, and non-deterministic -- relocating, not removing, the validation burden \citep{Ding2026Attribution}.

In contrast, human and deployment-grounded assessment recast the comparison in automated evaluators and human validation by changing which assumptions stay fixed relative to automated and model-based evaluators. The findings indicate that current RAG applications are mostly limited to domain-specific QA tasks, with systems still in prototype stages; industry requirements focus primarily on data protection, security, and quality, while issues such as ethics, bias, and scalability receive less attention; data preprocessing remains a key challenge, and system evaluation is predominantly conducted by humans rather than automated methods \citep{Brehme2025Retrieval}. Trade-off evidence from benchmark tasks does not travel cleanly across automated evaluators and human validation without matching budget and cost model conditions \citep{Caruzzo2026Deceptive}.

Across automated evaluators and human validation, the real split between automated and model-based evaluators and human and deployment-grounded assessment lies in which assumptions each route keeps fixed, not in one isolated result \citep{SaadFalcon2023Ares,Murugaraj2025Ragvue,Brehme2025Retrieval,Caruzzo2026Deceptive}. What remains uncertain in automated evaluators and human validation is not whether results exist, but whether they stay comparable once the setup shifts. Progress will remain local until papers report task, metric, horizon, and deployment conditions as one explicit package \citep{Murugaraj2025Ragvue,SaadFalcon2023Ares}.
\section{Robustness, Transfer \& Operational Validity}


The core issue in Robustness, Transfer \& Operational Validity is not breadth alone, but how robustness under shift and operational validity reshape what counts as a fair comparison \citep{Muthyam2026Systems,Feng2026Rethinking,Sun2026Robust,Li2026Mrag}. Sensitivity, adaptation, and query or corpus shift and Domain transfer, trustworthiness, and deployment each expose that pressure from a different angle, showing why robustness under shift and operational validity cannot be treated as interchangeable background choices \citep{Su2025Dynamic,Bang2026Feedback,Cahoon2025Optimizing,Ding2026Attribution}.
\subsection{Sensitivity, adaptation, and query or corpus shift}


The useful decision point in sensitivity, adaptation, and query or corpus shift is which commitments survive a change in benchmark, deployment setting, or reporting frame. Robustness evaluation should measure variance under query, corpus, and retrieval-policy shifts rather than only average answer accuracy \citep{Sun2026Robust,Muthyam2026Systems,Li2026Mrag,Bang2026Feedback}. This is the central pressure point: average benchmark performance can conceal instability when queries, corpora, retrieval policies, or feedback conditions change \citep{Sun2026Robust,Muthyam2026Systems,Li2026Mrag,Bang2026Feedback}. Trade-off evidence from benchmark tasks does not travel cleanly across sensitivity, adaptation, and query or corpus shift without matching budget and cost model conditions \citep{Sun2026Robust}.

Within sensitivity, adaptation, and query or corpus shift, papers grouped under perturbation and stability studies treat the underlying assumptions as part of the method rather than as reporting detail. Retrieval-Augmented Generation (RAG) systems are often evaluated using final answer accuracy, even though their failures can originate from preprocessing, retrieval, context packing, or generation \citep{Muthyam2026Systems}. Semantically equivalent queries with different syntactic forms may lead to different retrieval results, while irrelevant or misleading documents can further induce hallucinated answers \citep{Sun2026Robust}.

For perturbation and stability studies, reported gains in sensitivity, adaptation, and query or corpus shift are only informative when the evaluation setup keeps the relevant assumptions explicit. A narrower reading is warranted here because existing multi-path reasoning methods improve robustness by sampling multiple candidate answers and applying voting- or confidence-based selection, but they still face two limitations: diversity is often injected through uncontrollable decoding randomness, and answer evaluation is usually confined to a single query-induced evidence view \citep{Sun2026Robust}.

In contrast, adaptation and shift studies recast the comparison in sensitivity, adaptation, and query or corpus shift by changing which assumptions stay fixed relative to perturbation and stability studies. Retrieval-Augmented Generation (RAG) has become a foundational paradigm for equipping large language models (LLMs) with external knowledge, playing a critical role in information retrieval and knowledge-intensive applications \citep{Su2025Dynamic}. Existing evaluation protocols focus on overall accuracy and fail to capture how systems adapt after feedback is introduced \citep{Bang2026Feedback}.

The broader contrast in sensitivity, adaptation, and query or corpus shift only becomes clear once the comparison stays attached to the assumptions each route is actually holding fixed \citep{Sun2026Robust,Muthyam2026Systems,Li2026Mrag}. The hardest part of reading sensitivity, adaptation, and query or corpus shift is deciding which results survive a like-for-like comparison. A firmer comparison will depend on reporting task, metric, and deployment assumptions tightly enough to read perturbation and stability studies against adaptation and shift studies on the same footing \citep{Muthyam2026Systems,Sun2026Robust}.
\subsection{Domain transfer, trustworthiness, and deployment}


What most separates domain transfer, trustworthiness, and deployment is the split between trust and evaluator-transfer studies and domain and deployment studies, because their results move under different assumptions. Deployment claims require evidence that metrics and evaluators retain their meaning across domains, datasets, costs, and governance constraints \citep{Ding2026Attribution,Zhou2024Trustworthiness,Cahoon2025Optimizing,Kovalskii2026Domain}. The comparison becomes unstable at this point: evaluator agreement observed in one dataset or domain may not retain the same meaning after transfer to a different operating context \citep{Ding2026Attribution,Zhou2024Trustworthiness,Cahoon2025Optimizing,Kovalskii2026Domain}. In high-stakes information domains such as healthcare, where large language models (LLMs) can produce hallucinations or misinformation, retrieval-augmented generation (RAG) has been proposed as a mitigation strategy, grounding model outputs in external, domain-specific documents \citep{Javadi2025Contradictions}.

Papers grouped under trust and evaluator-transfer studies sharpen the contrast in domain transfer, trustworthiness, and deployment because changing those assumptions changes what the method is actually built to do. Although existing research mainly emphasizes accuracy and efficiency, the trustworthiness of RAG systems remains insufficiently explored \citep{Zhou2024Trustworthiness}. The evidence for trust and evaluator-transfer studies in domain transfer, trustworthiness, and deployment depends less on the headline number than on whether those assumptions remain comparable across papers. A more limited reading is the sounder one here because a prompt-based LLM judge avoids the chance-level collapses the automatic scorers suffer (no LFQA collapse) but is not uniformly best, \textasciitilde{}100x costlier, and non-deterministic -- relocating, not removing, the validation burden \citep{Ding2026Attribution}.

In contrast, seen through domain and deployment studies, the disagreement moves away from trust and evaluator-transfer studies and toward a different set of assumptions. Benchmarking across four diverse datasets highlights the strengths of different RAG methodologies, demonstrates TREX's ability to handle multiple open-domain QA types, and reveals the limitations of current evaluation methods \citep{Cahoon2025Optimizing}. When applied to heterogeneous corpora and multi-step queries, Naive RAG pipelines often degrade in quality due to flat knowledge representations and the absence of explicit workflows \citep{Kovalskii2026Domain}.

Across domain transfer, trustworthiness, and deployment, the real split between trust and evaluator-transfer studies and domain and deployment studies lies in which assumptions each route keeps fixed, not in one isolated result \citep{Ding2026Attribution,Zhou2024Trustworthiness,Cahoon2025Optimizing,Kovalskii2026Domain}. What remains uncertain in domain transfer, trustworthiness, and deployment is not whether results exist, but whether they stay comparable once the setup shifts. This instability has a concrete decision cost: a naive "best-on-average" rule for choosing an evaluator fails leave-one-dataset-out (mean held-out regret 0.172 AUROC, worse than fixing one scorer), so metric choice must be validated on the target dataset rather than learned from others \citep{Ding2026Attribution}. What is still missing is reporting that keeps task, metric, and deployment details visible when trust and evaluator-transfer studies and domain and deployment studies are compared \citep{Ding2026Attribution,Zhou2024Trustworthiness}.
\section{Discussion}


Across current approaches, the strongest pattern is not a single winning method but a recurring dependence on how design choices and evaluation protocols interact \citep{Sun2026Robust,Kim2024Learning,Salemi2024Evaluating,Taiwo2026Evaluating}. That interaction explains why apparently similar methods can report different outcomes under different assumptions, resource budgets, or evaluation setups \citep{Kovalskii2026Domain,Tian2026Predicting,Caruzzo2026Deceptive,Ding2026Attribution}. It also suggests that future work should treat evaluation design, protocol transparency, and practical constraints as first-class comparative variables rather than as appendix details \citep{Murugaraj2025Ragvue,Nematov2025Source,SaadFalcon2023Ares,Upadhyay2026Overview}.
\section{Conclusion}


Across this literature, performance is best read as evidence-bounded behavior shaped by coordinated choices about methodology, evaluation, and deployment setting \citep{Yao2026Authoritybench,An2026Fresco,Li2026Mrag,Guder2026Rage}. A survey that compares approaches without making protocol assumptions explicit will overstate some gains and understate some risks \citep{Brehme2025Retrieval,Packowski2024Optimizing,Muthyam2026Systems,Feng2026Rethinking}. By structuring the literature around shared lenses and comparable evidence, this comparative framing aims to make future surveys more reproducible, more stable, and easier to extend \citep{Su2025Dynamic,Bang2026Feedback,Cahoon2025Optimizing,Javadi2025Contradictions}.
\appendix

\section{Tables}



\begin{table}[H]
\centering
\scriptsize
\setlength{\tabcolsep}{2.5pt}
\renewcommand{\arraystretch}{1.0}
\caption{Cross-section comparison map for the survey areas.}
\label{tab:a1}
\begin{tabularx}{\linewidth}{YYYl}
\toprule
Survey area & Comparison focus & Representative evidence & Key refs \\
\midrule
3.1 Retrieval quality and evidence coverage & retrieval recall and rank quality, evidence sufficiency and context coverage & Retrieval diagnostic studies vs Retrieval utility and coverage on retrieval recall and rank quality & \citep{Salemi2024Evaluating,Taiwo2026Evaluating,Sun2026Robust,Kim2024Learning} \\
3.2 Generation faithfulness and source attribution & claim-level support and answer faithfulness, source-attribution correctness & Attribution and grounding vs Automated faithfulness evaluators on claim-level support and answer faithfulness & \citep{Ding2026Attribution,Nematov2025Source,SaadFalcon2023Ares,Murugaraj2025Ragvue} \\
4.1 Dataset, task, and metric design & corpus and query provenance, task coverage and difficulty & Evaluation frameworks and shared vs Domain and stress-test benchmarks on corpus and query provenance & \citep{Guder2026Rage,SaadFalcon2023Ares,Yao2026Authoritybench,Li2026Mrag} \\
4.2 Automated evaluators and human validation & rule-based versus learned versus LLM-based judging, agreement with human assessment & Automated and model-based evaluators vs Human and deployment-grounded on rule-based versus learned versus LLM-based judging & \citep{SaadFalcon2023Ares,Murugaraj2025Ragvue,Brehme2025Retrieval,Caruzzo2026Deceptive} \\
5.1 Sensitivity, adaptation, and query or corpus shift & query and corpus perturbation, retriever, chunker, and reranker sensitivity & Perturbation and stability studies vs Adaptation and shift studies on query and corpus perturbation & \citep{Sun2026Robust,Muthyam2026Systems,Li2026Mrag,Bang2026Feedback} \\
5.2 Domain transfer, trustworthiness, and deployment & cross-domain metric and evaluator transfer, fairness and contradiction handling & Trust and evaluator-transfer studies vs Domain and deployment studies on cross-domain metric and evaluator transfer & \citep{Ding2026Attribution,Zhou2024Trustworthiness,Cahoon2025Optimizing,Kovalskii2026Domain} \\
\bottomrule
\end{tabularx}
\end{table}

{\small
\bibliographystyle{plainnat}
\bibliography{../citations/ref}
}

\end{document}
